\chapter{Fazit}
In diesem Versuch wurden die Gammastrahlen-Energiespektren der drei untersuchten Elemente ($^{152}$Eu, $^{60}$Co und $^{137}$Cs) sowie von Gesteinsproben, die  mitgebracht wurden, aufgezeichnet und verglichen. Die Messungen erfolgten mit zwei verschiedenen Detektoren: einem NaI(Tl)-basierten Szintillationsdetektor (Sz-Detektor) und einem hochreinen Germaniumdetektor (HPGe-Detektor).

Die Energiekalibration zeigte bei beiden Detektoren einen linearen Zusammenhang zwischen Kanalnummer und Photonenenergie, sodass eine lineare Eichfunktion ausreichend war. Die Analyse der Halbwertsbreiten bestätigte die erwarteten Detektoreigenschaften: Der HPGe-Detektor erreichte mit FWHM-Werten zwischen $1,1\,\mathrm{keV}$ und $2,1\,\mathrm{keV}$ eine sehr hohe Energieauflösung,
während der NaI(Tl)-Detektor aufgrund der schlechteren Lichtstatistik deutlich breitere Peaks im Bereich von $10$ bis $75\,\mathrm{keV}$ zeigte.

Das absolute Effizienzverhältnis wurde ebenfalls bestimmt. Die absolute Effizienz des HPGe-Detektors betrug bei $661\,\mathrm{keV}$ etwa $(4,705 \pm0,042)\%$, während der Szintillationsdetektor aufgrund seiner höheren effektiven Ordnungszahl eine höhere Effizienz aufwies ($\epsilon(661\,\text{keV}) = (11,269\pm0,461)\%$). Dies lässt sich auf Unterschiede in der Ordnungszahl, dem Photonenfluss und der Geometrie der aktiven Teilchen zurückführen. Die Form der Effizienzkurve stimmt qualitativ mit den erwarteten Beiträgen des photoelektrischen Effekts und des Compton-Effekts überein.

%Darüber hinaus wurde die Leistungsfähigkeit der beiden Verbindungen ermittelt. Der HPGe-Detektor liefert eine Leistung von
%\begin{align*}
%    \mu = (\pm)\%.
%\end{align*}

Die Effizienzkurve wurde verwendet, um die Aktivität der beiden $^{60}$Co-Peaks zu bestimmen. Für den HPGe-Detektor ergeben sich daraus für die $^{60}$Co-Quelle:


\[
B_{\mathrm{HPGe}}(\text{Co}) = (26,1 \pm 0,9)\,\mathrm{kBq},
\]

während mit dem NaI(Tl)-Detektor

\[
B_{\mathrm{NaI}}(\text{Co}) = (8,5 \pm 0,6)\,\mathrm{kBq}
\]
Diese Werte liegen jedoch deutlich unter der durch das Zerfallsgesetz erwarteten Aktivität von $B_{\mathrm{Co}} = 36,1\,\mathrm{kBq}$. Die Abweichung ist auf systematische Unsicherheiten zurückzuführen die begrenzte Genauigkeit der Effizienzkurve sowie geometrische Unsicherheiten.

Die Untersuchung der Blumenerdeprobe mit dem HPGe-Detektor ermöglichte den Nachweis mehrerer charakteristischer $\gamma$-Linien, die den natürlichen Zerfallsreihen von Uran und Thorium sowie dem natürlich vorkommenden Isotop $^{40}$K zugeordnet werden konnten. Die gemessenen Energien stimmen innerhalb der Messunsicherheiten sehr gut mit den entsprechenden Literaturwerten überein, wodurch sowohl die Energiekalibration als auch die Zuordnung der beobachteten Nuklide bestätigt werden.

Die hohe Energieauflösung des HPGe-Detektors ermöglichte eine eindeutige Trennung der einzelnen Photopeaks und damit eine zuverlässige Identifikation der enthaltenen Radionuklide. Da jedoch kein separates Hintergrundspektrum zur Verfügung stand, konnte der Einfluss der natürlichen Umgebungsstrahlung nicht quantitativ berücksichtigt werden. Aus diesem Grund beschränkt sich die Auswertung auf den qualitativen Nachweis der in der Probe enthaltenen radioaktiven Nuklide; eine zuverlässige Bestimmung ihrer Aktivitäten war im Rahmen dieses Versuches nicht möglich.
